\documentclass[11pt,a4paper]{article}

\usepackage[margin=2.5cm]{geometry}
\usepackage[T1]{fontenc}
\usepackage[utf8]{inputenc}
\usepackage{lmodern}
\usepackage{microtype}
\usepackage{graphicx}
\usepackage{booktabs}
\usepackage{xcolor}
\usepackage[hidelinks]{hyperref}
\usepackage{xurl}
\usepackage[nameinlink]{cleveref}

\graphicspath{{figures/}}

\title{Urban Outfinders\\[0.3em]\large Part 2: Indexing and Evaluation}
\author{Marc Mallol\and Èric Dürr\and Oriol Pont}
\date{Information Retrieval and Web Analytics, 2026}

\begin{document}
\maketitle

\section{Introduction}

% What this part adds to Part 1: completed brand names, an inverted index,
% conjunctive (AND) matching, TF-IDF ranking and its evaluation. Outline of the sections. Where the code is:
% https://github.com/uripont/urban-outfinders, notebooks/part_2.ipynb.

\section{Indexing}

The code is in the \path{fashion_search} package: \path{index.py} builds the
inverted index and matches queries, and \path{ranking.py} ranks the matches.
Section~2 of \path{notebooks/part_2.ipynb} runs them.

\subsection{Inverted index}

\paragraph{Brand names completed.} Part~1 found that most brand names are cut
off in the source data (\emph{Pu}, \emph{True Bl}, \emph{REEB}), so a search for
\emph{puma} could not find a single Puma product. Product URLs start with the
full brand (\path{/puma-solid-men-round-neck-...}), so before indexing we
complete each brand whose letters run out inside a URL word with the rest of
that word: \emph{Pu} becomes \emph{Puma} and \emph{True Bl} becomes
\emph{True Blue}. This changes the brand of 26{,}798 of the 28{,}080 products.
Brands that do not match the start of their URL are left as they are. We also
strip the seller rating and notice glued to some seller names
(\emph{RetailNet4.5Seller changed. Check for\ldots} becomes \emph{RetailNet}).

\paragraph{Structure.} The index maps each term to the products that contain
it and how many times, $\{\mathit{term}: \{\mathit{pid}: \mathit{count}\}\}$.
The counts are what TF-IDF needs; the set of pids is the posting list of the
practical labs. Following Part~1, each product is indexed on one text that
merges its title, description, product details and brand, all processed by
the same \texttt{build\_terms} as the queries. Category and sub-category are
left out: they never contain query words and work better as filters.
\Cref{tab:index} gives the size of the result. Its vocabulary is very uneven:
half the terms appear in at most 4 products, while the most common
(\emph{wear}, \emph{wash}, \emph{cotton}) appear in about 70\% of them, mostly
through the product details.

\begin{table}[h]
  \centering
  \caption{Size of the inverted index.}
  \label{tab:index}
  \begin{tabular}{lr}
    \toprule
    Products & 28{,}080 \\
    Terms & 6{,}884 \\
    Postings (term, product pairs) & 905{,}629 \\
    Median products per term & 4 \\
    Terms in a single product & 2{,}043 \\
    \bottomrule
  \end{tabular}
\end{table}

\paragraph{Conjunctive matching.} A query matches the products that contain
every one of its terms. We intersect the posting lists of the query terms,
starting from the shortest, so each step can only shrink the set.
\Cref{tab:and} shows how each term narrows the validation queries: common terms
such as \emph{cotton} remove few products, while the garment words remove most.

\begin{table}[h]
  \centering
  \caption{Products matching each term, and all terms so far, for the two
  validation queries (terms after pre-processing).}
  \label{tab:and}
  \begin{tabular}{llrr}
    \toprule
    Query & Term & With the term & With all terms so far \\
    \midrule
    1 & women & 13{,}807 & 13{,}807 \\
      & full & 7{,}171 & 3{,}537 \\
      & sleev & 18{,}031 & 3{,}479 \\
      & sweatshirt & 1{,}419 & 671 \\
      & cotton & 19{,}889 & 498 \\
    \midrule
    2 & men & 13{,}517 & 13{,}517 \\
      & slim & 5{,}654 & 2{,}861 \\
      & jean & 2{,}436 & 368 \\
      & blue & 7{,}696 & 225 \\
    \bottomrule
  \end{tabular}
\end{table}

\subsection{TF-IDF ranking}

Each term $t$ of a product or a query gets the weight
\[
  w_t = \left(1 + \log \mathit{tf}_t\right) \cdot \log \frac{N}{\mathit{df}_t},
\]
where $\mathit{tf}_t$ is how many times $t$ appears in the product (or query),
$\mathit{df}_t$ is how many products contain $t$ and $N = 28{,}080$. The
logarithm on $\mathit{tf}$ keeps a word repeated across the product details
from dominating, and the IDF factor gives rare words more weight than words
that most products share. A product's score is the cosine similarity between
its weight vector and the query's, so longer product texts are not favoured
just for having more words. Products with equal scores, often the same item
listed by different sellers, are ordered by pid so results are reproducible.

Query~1 matches 498 products and query~2 matches 225. Their top results are
the expected garments (\emph{Full Sleeve Solid Women Sweatshirt},
\emph{Slim Men Blue Jeans}), from several brands.

\paragraph{Runs to evaluate.} \path{validation_labels.csv} judges only 20
products per query, and none of them is among the top 10 results of either
query. Counting unjudged products as non-relevant would give every metric a
score of zero and say nothing about the ranking. As in the practical labs, we
evaluate how the engine ranks the judged products: each run is the list of
judged products the engine returns, in its order. For query~1 the engine
returns all 20 judged products, including the 13 relevant ones. For query~2 it
returns only 8 of them, 6 of the 10 relevant ones, because the other four lack
a query term (\cref{sec:limitations}).

\subsection{Test queries}

% The five queries, how they were chosen (term frequencies, TF-IDF, ...) and
% what each one tests.

\section{Evaluation}

\subsection{Metrics}

% Each metric (P@K, R@K, AP@K, F1@K, MAP, MRR, NDCG): what it measures, how
% they differ, and the value of K used. Code: evaluation.py.

\subsection{Validation queries}

% Numbers only, rounded to three decimals: one table for the two validation
% queries, no text. The runs are the ones described in Section 2.2.

\subsection{Test queries}

% How the relevance labels were assigned (which documents were judged, by whom,
% what counts as relevant), the results table, and the analysis of the results.
% Labels: data/test_labels.csv.

\section{Limitations and improvements}
\label{sec:limitations}

Section~4 of the notebook holds the evidence for each problem below.

\paragraph{AND matching is too strict.} A product is lost if it lacks a single
query term, however well it matches the rest. Four of the ten relevant products
for \emph{men slim jeans blue} are never returned: one is a \emph{Tapered Fit
Men Blue Jeans} without \emph{slim}, and three are women's jeans that the
labels consider relevant but that do not contain \emph{men}. Common terms make
this worse: \emph{cotton} has an IDF of 0.35, so it barely changes the ranking,
yet it still excludes every product that does not mention it.
\emph{Fix:} rank with a disjunctive (OR) query, so missing terms lower the score
instead of removing the product, or drop query terms whose IDF is very low
before matching.

\paragraph{Different words for the same thing.} Products and users name the
same garment differently, and a query finds only the products that use its
word: \emph{sweater} matches 149 products and \emph{sweatshirt} 1{,}419,
\emph{sneakers} 780 and \emph{shoes} 526, with little overlap. Compounds split
the same way: \emph{track pants} and \emph{trackpants} are different terms.
Part~1 unified a few such variants (\emph{tee}, \emph{t-shirt}); a fixed list
does not scale. \emph{Fix:} expand queries with synonyms from a hand-made
garment list or from word embeddings (Part~3), and join compounds that appear
both ways in the catalogue.

\paragraph{One merged field.} A term counts the same in the title as in a long
list of product details, although the title says what the product is and the
details mostly describe its fabric and care. \emph{Fix:} index the fields
separately and weight title matches more, as BM25F does.

\paragraph{Near-duplicates and no quality signal.} Many products share a title
and differ only in colour or seller, so the top 20 results of the validation
queries contain only 3 and 4 distinct titles. TF-IDF also ignores rating,
price and stock: four of the top 20 results of query~1 are out of stock.
\emph{Fix:} combine the text score with rating, discount and availability (the
custom score of Part~3), and group or diversify near-identical products.

\paragraph{Incomplete relevance judgements.} The validation labels judge 20
products per query, while the queries match 498 and 225 products, so the
metrics only measure how the judged products are ordered, not the quality of
what users would see first. \emph{Fix:} judge the top results of every system
being compared (pooling).

\section{Conclusions}

\section*{Use of AI tools}

% Check this with the group before submitting.

\end{document}
